% Source: https://artofproblemsolving.com/community/c3h3269739_how_to_find_number_of_digits_in_a_really_large_number
% Stable identity: aops-topic-3269739-q01
A positive integer $N$ has $d$ digits exactly when $10^{d-1}\le N<10^d$, so $\boxed{d=\lfloor y\log_{10}x\rfloor+1}$. For $2^{100}$, $100\log_{10}2=30.102999\ldots$, giving $31$ digits. Its leading digit is $\lfloor10^{0.102999\ldots}\rfloor=1$, and the four-cycle of powers of two gives last digit $6$. Their sum is $\boxed7$.
